% GENERATED FILE. Edit canonical lessons, then run npm run generate:books.
% canonical-source-hash: fnv1a64:08e6d9461b9e6dbd
% canonical-lessons: HI-C284-khush, HI-C284-darlagta, HI-C284-bannachahta, HI-R284-feelings-and-hopes

\chapter{Feelings and Hopes}
\label{ch:hi-feelings-and-hopes}
\hlchaptermodality{284}

\begin{chapteropening}
\textbf{By the end of this chapter:} \emph{I can say how I feel (mujhe ḍar lagtā hai, maiṁ khuś hūṁ, mujhe acchā lagā), what I want to become and what I hope for (maiṁ adhyāpak bannā cāhtā hūṁ, mujhe ummīd hai ki …).}

Complete HI-R284-feelings-and-hopes: say how you feel, what you want to become and what you hope for.
\end{chapteropening}

\section[\texorpdfstring{\hi{ख़ुश} (khuś)}{ख़ुश (khuś)}]{\hi{ख़ुश} (khuś) — happy, glad}

\label{lesson:HI-C284-khush}



\begin{quote}
Before the new one: say the line from your first meeting. (\textbf{\hi{आपसे} \hi{मिलकर} \hi{ख़ुशी} \hi{हुई}}.) Its \textbf{\hi{ख़ुशी}} is \textquotedblleft{}happiness\textquotedblright{}. Then \textquotedblleft{}sad\textquotedblright{}. (\textbf{\hi{उदास}}, \emph{udās}.)
\end{quote}

\subsection*{You'll want to know: \hi{ख़ुश}}
\textbf{\hi{ख़ुश}} — \emph{khuś} — \textquotedblleft{}happy, glad\textquotedblright{}.

\begin{quote}
\textbf{\hi{आज} \hi{मैं} \hi{बहुत} \hi{ख़ुश} \hi{हूँ।}}
\end{quote}

\begin{quote}
\emph{āj maiṁ bahut khuś hūṁ.} — \textquotedblleft{}I am very happy today.\textquotedblright{}
\end{quote}

\textbf{\hi{ख़ुश}} never changes: \textbf{\hi{वह} \hi{ख़ुश} \hi{है}} for him and for her, \textbf{\hi{हम} \hi{ख़ुश} \hi{हैं}} for us. It is the word \textbf{\hi{ख़ुशी}} is built on: \textbf{\hi{ख़ुश}}, \textquotedblleft{}happy\textquotedblright{}, and \textbf{\hi{ख़ुशी}} (\emph{khuśī}), \textquotedblleft{}happiness\textquotedblright{}.

With \textbf{\hi{होना}} in the past, it says the gladness came: \textbf{\hi{वह} \hi{बहुत} \hi{ख़ुश} \hi{हुई}} (\emph{vah bahut khuś huī}), \textquotedblleft{}she was very pleased\textquotedblright{}. The one you are pleased with takes \textbf{\hi{से}}: \textbf{\hi{मैं} \hi{तुमसे} \hi{ख़ुश} \hi{हूँ}} (\emph{maiṁ tumse khuś hūṁ}), \textquotedblleft{}I am pleased with you\textquotedblright{}.

\begin{cousinweb}
\textbf{\hi{ख़ुश}} (\emph{khuś}) is usually traced to Persian \emph{khush}, \textquotedblleft{}happy, pleasant\textquotedblright{}. \textbf{\hi{ख़ुशी}} adds \textbf{-\hi{ई}}, the ending of many Hindi nouns of feeling, and \textbf{\hi{ख़ुशबू}}, \textquotedblleft{}fragrance\textquotedblright{}, opens with the same word.
\end{cousinweb}

\subsection*{Your turn}
Say these aloud:

\begin{itemize}
\raggedright
  \item \emph{khuś}, then \emph{khuśī}
  \item \textquotedblleft{}I am very happy today\textquotedblright{} — \emph{āj maiṁ bahut khuś hūṁ}
  \item \textquotedblleft{}She was very pleased\textquotedblright{} — \emph{vah bahut khuś huī}
\end{itemize}

\begin{tcolorbox}[breakable,colback=teal!4,colframe=teal!35!black,title={Before you move on}]
What does \textbf{\hi{ख़ुश}} mean? (\textquotedblleft{}Happy, glad\textquotedblright{}.) Does it change for a woman? (No: \textbf{\hi{वह} \hi{ख़ुश} \hi{है}}.) \textquotedblleft{}Happiness\textquotedblright{}? (\textbf{\hi{ख़ुशी}}.)
\end{tcolorbox}
\section[\texorpdfstring{\hi{मुझे} \hi{डर} \hi{लगता} \hi{है} (mujhe ḍar lagtā hai)}{मुझे डर लगता है (mujhe ḍar lagtā hai)}]{\hi{मुझे} \hi{डर} \hi{लगता} \hi{है} (mujhe ḍar lagtā hai) — saying how you feel}

\label{lesson:HI-C284-darlagta}



\begin{quote}
Say \textquotedblleft{}I'm hungry\textquotedblright{} and \textquotedblleft{}I am happy\textquotedblright{}. (\emph{mujhe bhūkh lagī hai}, \emph{maiṁ khuś hūṁ}.) Two ways of feeling, and this lesson sets them side by side.
\end{quote}

\subsection*{You'll want to know: a feeling that comes to you}
\begin{quote}
\textbf{\hi{मुझे} \hi{कुत्तों} \hi{से} \hi{डर} \hi{लगता} \hi{है।}}
\end{quote}

\begin{quote}
\emph{mujhe kuttoṁ se ḍar lagtā hai.} — \textquotedblleft{}I'm afraid of dogs.\textquotedblright{}
\end{quote}

Some feelings come to you, and Hindi says so. The one who feels takes \textbf{\hi{मुझे}}, the feeling is the subject of \textbf{\hi{लगना}}, and what you fear takes \textbf{\hi{से}}.

\noindent
\begin{tabularx}{\linewidth}{@{}>{\raggedright\arraybackslash}X>{\raggedright\arraybackslash}X@{}}
\toprule
\textbf{the feeling} & \textbf{the sentence} \\
\midrule
\textbf{\hi{डर}}, masculine & \textbf{\hi{मुझे} \hi{डर} \hi{लगता} \hi{है}} (\emph{mujhe ḍar lagtā hai}) — I'm afraid \\
\textbf{\hi{ठंड}}, feminine & \textbf{\hi{मुझे} \hi{ठंड} \hi{लग} \hi{रही} \hi{है}} (\emph{mujhe ṭhaṇḍ lag rahī hai}) — I'm cold \\
how it seemed & \textbf{\hi{मुझे} \hi{अच्छा} \hi{लगा}} (\emph{mujhe acchā lagā}) — I liked it; it felt good \\
\bottomrule
\end{tabularx}

Anger comes with \textbf{\hi{आना}}, \textquotedblleft{}to come\textquotedblright{}: \textbf{\hi{मुझे} \hi{गुस्सा} \hi{आया}} (\emph{mujhe gussā āyā}), \textquotedblleft{}I got angry\textquotedblright{}.

\subsection*{You'll want to know: a feeling you are in}
\begin{quote}
\textbf{\hi{मीरा} \hi{आज} \hi{उदास} \hi{है।}}
\end{quote}

\begin{quote}
\emph{Mīrā āj udās hai.} — \textquotedblleft{}Mira is sad today.\textquotedblright{}
\end{quote}

Other feelings are a state you are in, said with \textbf{\hi{होना}}: \textbf{\hi{मैं} \hi{ख़ुश} \hi{हूँ}}; \textbf{\hi{वह} \hi{परेशान} \hi{है}} (\emph{vah pareśān hai}), \textquotedblleft{}he is worried\textquotedblright{}. \textbf{\hi{ख़ुश}}, \textbf{\hi{उदास}} and \textbf{\hi{परेशान}} never change; \textbf{\hi{अकेला}} agrees: \textbf{\hi{वह} \hi{अकेली} \hi{है}} (\emph{vah akelī hai}), \textquotedblleft{}she is alone\textquotedblright{}.

With \textbf{\hi{हो} \hi{गया}}, the feeling sets in: \textbf{\hi{वह} \hi{उदास} \hi{हो} \hi{गई}} (\emph{vah udās ho gaī}), \textquotedblleft{}she became sad\textquotedblright{}. To ask: \textbf{\hi{क्या} \hi{हुआ}? \hi{तुम} \hi{परेशान} \hi{क्यों} \hi{हो}?} (\emph{kyā huā? tum pareśān kyoṁ ho?}), \textquotedblleft{}What happened? Why are you worried?\textquotedblright{}

\subsection*{Your turn}
Say these aloud:

\begin{itemize}
\raggedright
  \item \textquotedblleft{}I'm cold\textquotedblright{} — \emph{mujhe ṭhaṇḍ lag rahī hai}
  \item \textquotedblleft{}I liked the food\textquotedblright{} — \emph{mujhe khānā acchā lagā}
  \item \textquotedblleft{}Why are you worried?\textquotedblright{} — \emph{tum pareśān kyoṁ ho?}
\end{itemize}

\begin{tcolorbox}[breakable,colback=teal!4,colframe=teal!35!black,title={Before you move on}]
\textquotedblleft{}I'm afraid\textquotedblright{}? (\textbf{\hi{मुझे} \hi{डर} \hi{लगता} \hi{है}}.) \textquotedblleft{}I got angry\textquotedblright{}? (\textbf{\hi{मुझे} \hi{गुस्सा} \hi{आया}}.) \textquotedblleft{}She is sad\textquotedblright{}? (\textbf{\hi{वह} \hi{उदास} \hi{है}}.) \textquotedblleft{}She became sad\textquotedblright{}? (\textbf{\hi{वह} \hi{उदास} \hi{हो} \hi{गई}}.) Which feelings take \textbf{\hi{मुझे}}? (The ones that come to you: \textbf{\hi{डर}}, \textbf{\hi{ठंड}}, \textbf{\hi{गुस्सा}}.)
\end{tcolorbox}
\section[\texorpdfstring{\hi{मैं} \hi{अध्यापक} \hi{बनना} … (maiṁ adhyāpak …)}{मैं अध्यापक बनना … (maiṁ adhyāpak …)}]{\hi{मैं} \hi{अध्यापक} \hi{बनना} \hi{चाहता} \hi{हूँ} (maiṁ adhyāpak bannā cāhtā hūṁ) — what I want to become}

\label{lesson:HI-C284-bannachahta}



\begin{quote}
Say \textquotedblleft{}I want to go\textquotedblright{}. (\textbf{\hi{मैं} \hi{जाना} \hi{चाहता} \hi{हूँ}}.) Then \textbf{\hi{बनना}}, \textquotedblleft{}to become\textquotedblright{}, and \textbf{\hi{सपना}}, \textquotedblleft{}a dream\textquotedblright{}.
\end{quote}

\subsection*{You'll want to know: wanting to become}
\begin{quote}
\textbf{\hi{मैं} \hi{अध्यापक} \hi{बनना} \hi{चाहता} \hi{हूँ।}}
\end{quote}

\begin{quote}
\emph{maiṁ adhyāpak bannā cāhtā hūṁ.} — \textquotedblleft{}I want to become a teacher.\textquotedblright{}
\end{quote}

\textbf{\hi{चाहना}} takes a verb in \textbf{-\hi{ना}}, as in \textbf{\hi{मैं} \hi{जाना} \hi{चाहता} \hi{हूँ}}. With \textbf{\hi{बनना}}, \textquotedblleft{}to become\textquotedblright{}, it says what you want to be. A woman says \textbf{\hi{चाहती} \hi{हूँ}}: \textbf{\hi{मैं} \hi{डाक्टर} \hi{बनना} \hi{चाहती} \hi{हूँ}} (\emph{maiṁ ḍākṭar bannā cāhtī hūṁ}).

\textquotedblleft{}When I grow up\textquotedblright{} is \textbf{\hi{बड़ा} \hi{होकर}} (\emph{baṛā hokar}), the \textbf{-\hi{कर}} of \textbf{\hi{होना}}: \textbf{\hi{बड़ा} \hi{होकर} \hi{मैं} \hi{अध्यापक} \hi{बनूँगा}} (\emph{baṛā hokar maiṁ adhyāpak banūṅgā}), \textquotedblleft{}when I grow up I will be a teacher\textquotedblright{}. A girl says \textbf{\hi{बड़ी} \hi{होकर} … \hi{बनूँगी}}.

To ask: \textbf{\hi{तुम} \hi{क्या} \hi{बनना} \hi{चाहते} \hi{हो}?} (\emph{tum kyā bannā cāhte ho?}), \textquotedblleft{}What do you want to be?\textquotedblright{}

\subsection*{You'll want to know: hopes and dreams}
\begin{quote}
\textbf{\hi{मेरा} \hi{सपना} \hi{है} \hi{कि} \hi{एक} \hi{दिन} \hi{मेरा} \hi{अपना} \hi{घर} \hi{होगा।}}
\end{quote}

\begin{quote}
\emph{merā sapnā hai ki ek din merā apnā ghar hogā.} — \textquotedblleft{}My dream is that one day I will have a house of my own.\textquotedblright{}
\end{quote}

\textbf{\hi{सपना}} is masculine: \textbf{\hi{मेरा} \hi{सपना}}. After \textbf{\hi{कि}}, the dream is told in the future, with \textbf{\hi{होगा}}. A dream can also be a verb in \textbf{-\hi{ना}}: \textbf{\hi{मेरा} \hi{सपना} \hi{अध्यापक} \hi{बनना} \hi{है}} (\emph{merā sapnā adhyāpak bannā hai}), \textquotedblleft{}my dream is to become a teacher\textquotedblright{}.

Hope works the same way: \textbf{\hi{मुझे} \hi{उम्मीद} \hi{है} \hi{कि} …} (\emph{mujhe ummīd hai ki …}), \textquotedblleft{}I hope that …\textquotedblright{}, with the future after it: \textbf{\hi{मुझे} \hi{उम्मीद} \hi{है} \hi{कि} \hi{तुम} \hi{आओगे}} (\emph{mujhe ummīd hai ki tum āoge}), \textquotedblleft{}I hope you will come\textquotedblright{}.

\subsection*{Your turn}
Say these aloud:

\begin{itemize}
\raggedright
  \item \textquotedblleft{}My dream is to become a doctor\textquotedblright{} — \emph{merā sapnā ḍākṭar bannā hai}
  \item \textquotedblleft{}I hope it won't rain tomorrow\textquotedblright{} — \emph{mujhe ummīd hai ki kal bāriś nahīṁ hogī}
\end{itemize}

\begin{tcolorbox}[breakable,colback=teal!4,colframe=teal!35!black,title={Before you move on}]
\textquotedblleft{}I want to become a teacher\textquotedblright{}? (\textbf{\hi{मैं} \hi{अध्यापक} \hi{बनना} \hi{चाहता} \hi{हूँ}}, or \textbf{\hi{चाहती} \hi{हूँ}}.) \textquotedblleft{}When I grow up\textquotedblright{}? (\textbf{\hi{बड़ा} \hi{होकर}}.) \textquotedblleft{}I hope that you will come\textquotedblright{}? (\textbf{\hi{मुझे} \hi{उम्मीद} \hi{है} \hi{कि} \hi{तुम} \hi{आओगे}}.) Why \textbf{\hi{मेरा} \hi{सपना}}? (\textbf{\hi{सपना}} is masculine.)
\end{tcolorbox}
\section[(review)]{Review — feelings and hopes}

\label{lesson:HI-R284-feelings-and-hopes}



\begin{quote}
Say \textquotedblleft{}I am happy\textquotedblright{} and \textquotedblleft{}I'm afraid\textquotedblright{}. (\emph{maiṁ khuś hūṁ}, \emph{mujhe ḍar lagtā hai}.)
\end{quote}

\subsection*{Your turn}
Say these aloud:

\begin{itemize}
\raggedright
  \item I'm afraid of dogs. I'm cold.
  \item I liked your village very much.
  \item Why are you sad? What happened?
  \item What do you want to become? I want to become a doctor.
  \item I hope that next year I will go to the village.
\end{itemize}

\begin{tcolorbox}[breakable,colback=teal!4,colframe=teal!35!black,title={Before you move on}]
\textquotedblleft{}I liked it\textquotedblright{}? (\textbf{\hi{मुझे} \hi{अच्छा} \hi{लगा}}.) \textquotedblleft{}I got angry\textquotedblright{}? (\textbf{\hi{मुझे} \hi{गुस्सा} \hi{आया}}.) \textquotedblleft{}She is happy\textquotedblright{}? (\textbf{\hi{वह} \hi{ख़ुश} \hi{है}}: \textbf{\hi{ख़ुश}} does not change.) \textquotedblleft{}I want to become …\textquotedblright{}? (\textbf{\hi{मैं} … \hi{बनना} \hi{चाहता} \hi{हूँ}}.) \textquotedblleft{}I hope that …\textquotedblright{}? (\textbf{\hi{मुझे} \hi{उम्मीद} \hi{है} \hi{कि} …}, then the future.)
\end{tcolorbox}
